% Problem record op_ae968c5086f34fea. This TeX file is derived from
% database/problems_json/op_ae968c5086f34fea.json by scripts/sync-tex.mjs; edit the JSON record
% and rerun the script rather than editing this file.
\section{Optimal asymptotic distillation yield of dephased T states}
\label{sec:op_ae968c5086f34fea}

\paragraph{Problem.}
What is the optimal catalyst-free asymptotic $T$-state yield of a dephased single-qubit $T$ state?

A stabilizer operation is constructed from stabilizer ancillas, Clifford unitaries, Pauli measurements, classical randomness and feedforward, and discarding systems. Heralding is permitted where stated. It does not mean an arbitrary stabilizer-preserving channel, and no magic catalyst is free. These distinctions affect conversion theorems \sourcecite{ref:ae96-zurel26}{Zurel26}, \sourcecite{ref:ae96-fang26}{Fang26}.

Throughout, the modern phase-gate magic state is

\begin{equation}
|T\rangle=\frac{|0\rangle+e^{i\pi/4}|1\rangle}{\sqrt2},
\qquad \tau=|T\rangle\langle T|.
\label{eq:ae96-1}
\end{equation}

Consider the one-parameter noise model

\begin{equation}
\rho_p=(1-p)\tau+pZ\tau Z,\qquad 0\leq p\leq1.
\label{eq:ae96-2}
\end{equation}

Define its catalyst-free asymptotic $T$-yield by

\begin{equation}
D_T(\rho_p)=\sup\left\{R:\ \exists\text{ stabilizer CPTP maps }\Lambda_n,
\quad
\frac12\left\|\Lambda_n(\rho_p^{\otimes n})-
\tau^{\otimes\lfloor Rn\rfloor}\right\|_1\longrightarrow0
\right\}.
\label{eq:ae96-3}
\end{equation}

Here the error concerns the joint output, not only individual marginals. A heralded implementation can be included by requiring its failure probability to vanish and specifying an output on failure.

Determine $D_T(\rho_p)$, with matching achievable and converse bounds, already for a fixed noise level such as $p=0.01$.

This is a precise specialization of the published open optimization question about asymptotically optimal distillation protocols \sourcecite{ref:ae96-wills25}{Wills25} (Open Question 4). It is not a separately named conjecture, and no particular candidate formula is asserted to be correct.

The displayed definitions, constraints, and target bounds are recorded in Eqs.~\eqref{eq:ae96-1}, \eqref{eq:ae96-2}, \eqref{eq:ae96-3}.

\subsection*{Status}

Unsolved

\subsection*{Source}

This precise formulation is editor wording based on the unresolved direction and limitations documented in the cited primary literature \sourcecite{ref:ae96-wills25}{Wills25}\sourcecite{ref:ae96-rubboli24}{Rubboli24}; it is not presented as a verbatim conjecture of those authors.

\subsection*{Progress}

\begin{itemize}
  \item Constant-overhead existence is solved. Wills, Hsieh, and Yamasaki constructed protocols with positive asymptotic yield for sufficiently good inputs; the paper appeared online in \emph{Nature Physics} on September 16, 2025. The open issue here is the largest yield for a specified input, not whether overhead exponent zero is possible \sourcecite{ref:ae96-wills25}{Wills25}. A fixed preliminary purification stage extends such constructions to inputs within an established distillation basin.

  \item A computable converse is available. Let

\begin{equation}
q=\cos^2(\pi/8)=\frac{1+1/\sqrt2}{2}.
\label{eq:ae96-4}
\end{equation}

The displayed definitions, constraints, and target bounds are recorded in Eqs.~\eqref{eq:ae96-4}.

  \item The stabilizer boundary on this axis occurs at $p_*=1-q$. Additivity and conversion bounds for relative entropy of magic give, for $0\leq p<p_*$,

\begin{equation}
D_T(\rho_p)\leq
\frac{(1-p)\log_2\frac{1-p}{q}
+p\log_2\frac{p}{1-q}}
{-\log_2q}.
\label{eq:ae96-5}
\end{equation}

The displayed definitions, constraints, and target bounds are recorded in Eqs.~\eqref{eq:ae96-5}.

  \item This specializes the results of Rubboli, Takagi, and Tomamichel to the commuting optimizer $q\tau+(1-q)Z\tau Z$ \sourcecite{ref:ae96-rubboli24}{Rubboli24} (Sections 4, 6 and 8). Direct substitution yields

\begin{equation}
D_T(\rho_{0.01})\leq0.757659\ldots.
\label{eq:ae96-6}
\end{equation}

The displayed definitions, constraints, and target bounds are recorded in Eqs.~\eqref{eq:ae96-6}.

  \item This number is a converse bound, not an achieved yield or a conjectured equality.

  \item May 28, 2026. Ehara and Takagi construct protocols with output counts arbitrarily close to linear in a sublinear sense, including $n^{1-\eta}$ for every fixed $\eta>0$. Their result concerns the relation between rate scaling and overhead exponents; it does not identify the optimal constant multiplying $n$ \sourcecite{ref:ae96-ehara26}{Ehara26}.

  \item Catalysis is a different task. Fang and Liu’s work, updated September 2, 2026, can reduce one-shot input count to one using a catalyst and reduced success probability. It neither supplies the missing catalyst-free optimal yield nor justifies ignoring failed attempts in rate accounting \sourcecite{ref:ae96-fang26}{Fang26} (Theorem 7).

  \item August 10, 2026. New binary-extension-field code constructions target practical multi-qubit magic-state protocols; they do not establish a matching optimum for the quantity above \sourcecite{ref:ae96-gong26}{Gong26}.

  \item Status: open quantitative optimization, despite the solved scaling-exponent problem.
\end{itemize}

\subsection*{References}

\begin{enumerate}[label={},leftmargin=4.8em,labelsep=0.5em]
  \footnotesize
  \item[\textup{[Zurel26]}]\label{ref:ae96-zurel26}
  M. Zurel, S. Jana, and N. de Silva, \emph{High-threshold magic state distillation with quantum quadratic residue codes}, \href{https://arxiv.org/html/2603.18560v1}{arXiv:2603.18560v1}, March 19, 2026. Introduction explicitly states the open conjectures; Sections 4.2–4.3 give the new constructions.

  \item[\textup{[Fang26]}]\label{ref:ae96-fang26}
  K. Fang and Z.-W. Liu, \emph{One-shot distillation with constant overhead using catalysts}, Nature Communications 17, 6010 (2026); \href{https://arxiv.org/html/2410.14547v3}{arXiv:2410.14547v3}, September 2, 2026. See Definitions 1–2 and Theorem 7.

  \item[\textup{[Wills25]}]\label{ref:ae96-wills25}
  A. Wills, M.-H. Hsieh, and H. Yamasaki, \emph{Constant-Overhead Magic State Distillation}, Nature Physics 21, 1842–1846 (2025); \href{https://arxiv.org/html/2408.07764v2}{arXiv:2408.07764}. The extended arXiv version explicitly asks for optimal overhead and asymptotically optimal protocols in Open Question 4.

  \item[\textup{[Rubboli24]}]\label{ref:ae96-rubboli24}
  R. Rubboli, R. Takagi, and M. Tomamichel, \emph{Mixed-state additivity properties of magic monotones based on quantum relative entropies for single-qubit states and beyond}, Quantum 8, 1492 (2024); \href{https://arxiv.org/html/2307.08258v3}{arXiv:2307.08258v3}.

  \item[\textup{[Ehara26]}]\label{ref:ae96-ehara26}
  K. Ehara and R. Takagi, \emph{Asymptotic magic state distillation with almost linear rate}, \href{https://arxiv.org/html/2605.30108v1}{arXiv:2605.30108v1}, May 28, 2026. See Equation (1), Theorem 1, and the success-probability analysis.

  \item[\textup{[Gong26]}]\label{ref:ae96-gong26}
  A. Gong, C. A. Pattison, P. Rall, and A. Wills, \emph{Magic State Distillation via Codes over Binary Extension Fields}, \href{https://arxiv.org/abs/2608.09727}{arXiv:2608.09727}, August 10, 2026.
\end{enumerate}

\subsection*{Comment}

This problem addresses an operationally meaningful constant that overhead exponents conceal. Two protocols can both have constant overhead yet consume very different numbers of noisy states for a given output volume.

A manageable starting point is to tighten either side for one fixed rational $p$, with every auxiliary magic state and failed attempt included in the accounting. An apparent formula derived under all stabilizer-preserving maps, nonzero-probability postselection, or borrowed magic should not be substituted for the catalyst-free, vanishing-global-error definition.

\subsection*{Field}

Quantum Resource Theory

\subsection*{Topic}

Quantum magic; Resource conversion; Additivity and regularization

\subsection*{ID}

\texttt{op\_ae968c5086f34fea}
